\section{A decision rule for spectral analysis}
\label{sec:book-spectral-decision-primer}

The word ``spectrum'' hides several different questions.  The operator type determines which
decomposition is informative.

For a symmetric matrix $L=L^{\mathsf T}$, the eigendecomposition
\begin{equation}
L=V\Lambda V^{\mathsf T}
\label{eq:book-primer-symmetric-spectrum}
\end{equation}
uses an orthonormal basis.  Its Rayleigh quotient
$u^{\mathsf T}Lu/\|u\|_2^2$ measures the quadratic energy per unit signal, so low Laplacian
eigenvalues identify smooth directions on the declared graph.

For a rectangular map, or for a directed map whose left and right spaces play different roles, use
the singular-value decomposition
\begin{equation}
O=U\Sigma V^{\mathsf T}.
\label{eq:book-primer-directed-svd}
\end{equation}
A right singular vector is an input direction, the matching left singular vector is its output
direction, and the singular value is the amplification.  A square map is nonnormal when
$OO^{\mathsf T}\ne O^{\mathsf T}O$.  Its eigenvectors can be nonorthogonal, so its SVD and its
eigendecomposition answer different questions.

A row-stochastic route $P$ adds a probabilistic question.  Its stationary law satisfies
\begin{equation}
\pi^{\mathsf T}P=\pi^{\mathsf T},
\qquad \sum_i\pi_i=1,
\label{eq:book-primer-stationary-law}
\end{equation}
and describes mass preserved by repeated routing.  Reversibility asks whether every stationary
edge flux balances its reverse:
\begin{equation}
\pi_iP_{ij}=\pi_jP_{ji}.
\label{eq:book-primer-detailed-balance}
\end{equation}
Entrywise symmetry is stronger than this flux condition.

For example, a lazy random walk on the three-node path has
\begin{equation}
P_{\mathrm{rev}}
=\begin{pmatrix}
1/2&1/2&0\\
1/4&1/2&1/4\\
0&1/2&1/2
\end{pmatrix},
\qquad
\pi=\begin{pmatrix}1/4\\1/2\\1/4\end{pmatrix}.
\label{eq:book-primer-reversible-asymmetric-walk}
\end{equation}
The entries $P_{12}=1/2$ and $P_{21}=1/4$ differ, while their stationary fluxes both equal $1/8$.
The route is asymmetric in Euclidean coordinates and reversible in the $\pi$-weighted geometry.

By contrast, the biased three-cycle
\begin{equation}
P_{\circlearrowright}
=\begin{pmatrix}
1/4&1/2&1/4\\
1/4&1/4&1/2\\
1/2&1/4&1/4
\end{pmatrix}
\label{eq:book-primer-directed-cycle}
\end{equation}
has the uniform stationary law because it is doubly stochastic.  Its clockwise flux exceeds its
counterclockwise flux on every edge, so detailed balance fails and a nonzero circulation remains.
With the entropy-production convention used later, this example has
$\operatorname{EP}(P_{\circlearrowright})=\tfrac14\log2>0$.

Finally, an individual eigenvector is not an invariant object inside a repeated or nearly repeated
spectral band.  If the columns of $V_B$ span such a band, another orthonormal basis has the form
$V_BQ$ with $Q^{\mathsf T}Q=I$, while
\begin{equation}
(V_BQ)(V_BQ)^{\mathsf T}=V_BV_B^{\mathsf T}.
\label{eq:book-primer-band-projector}
\end{equation}
The projector records the stable subspace even when the displayed eigenvectors rotate.  Thus the
minimal decision rule is: use eigenpairs for self-adjoint quadratic geometry, singular vectors for
directed input--output amplification, stationary flux for repeated stochastic routing, and
subspace projectors for clustered spectral modes.
